%
% crosshost_connection_lifecycle_no_handshake_gate_buggy.tex
%
% FIDELITY CONTROL (variant A) for crosshost_connection_lifecycle.tex
% (W12, lux-833p).  Removes ONLY the Gate 1 ordering: Accept makes a
% freshly-connected socket readable immediately, before its TLS handshake
% has completed (the "handshake-before-accept ordering removed" defect the
% design's fidelity requirement names).  Reproduces the defect: I1's
% negation (an identified-but-not-handshaken connection) is REACHABLE here
% (goal FOUND), whereas it is unreachable in the correct spec.  I2 still
% holds here -- this variant isolates the Gate 1 failure.
% The authoritative model is crosshost_connection_lifecycle.tex.
%
% Models the code in:
%   src/punt_lux/display/cross_host_listener.py    -- CrossHostListener
%       (Gate 1: the pending-handshake set; a socket is not in the reader
%        set until do_handshake() has returned "ready")
%   src/punt_lux/display/cross_host_verification.py -- CrossHostVerification
%       (Gate 2: the SAN/hub_id cross-check)
%   src/punt_lux/display/hub_reconciliation.py     -- handle_connect
%       (Gate 2: reject-or-register-identity)
%   src/punt_lux/display/identity_guard.py         -- reject_scene_unless_hub
%       (content requires an identified kind="hub" fd)
% and the invariant stated in
%   docs/architecture/system.tex "Cross-Host Transport and Trust",
%   subsection "Invariants", item 1, and subsection "Coexistence with the
%   Local Fast Path" ("Two distinct gates, not one").
%
% Type-check:  fuzz docs/crosshost_connection_lifecycle_no_handshake_gate_buggy.tex
% Model-check (goal predicates in the Verification section):
%   probcli docs/crosshost_connection_lifecycle_no_handshake_gate_buggy.tex \
%       -model_check -p DEFAULT_SETSIZE 3
%
% Formatting follows the fuzz house rules: \quad~ for continuation indent,
% schema lines under 80 columns, predicates broken at \land/\lor/\implies.
%

\documentclass[a4paper,10pt,fleqn]{article}
\usepackage[margin=1in]{geometry}
\usepackage{fuzz}
\usepackage[colorlinks=true,linkcolor=blue,citecolor=blue,urlcolor=blue]{hyperref}

\begin{document}

\title{Cross-Host Connection Lifecycle:
a Fidelity Control (variant A, no Gate 1, BUGGY)}
\author{W12 variant --- handshake-before-accept ordering removed, which
violates DES-090 Invariant 1}
\date{September 2026}
\maketitle

\noindent\fbox{\parbox{\dimexpr\textwidth-2\fboxsep-2\fboxrule}{%
\textbf{Fidelity control --- BUGGY by construction.} This variant makes a
freshly-accepted socket readable before its TLS handshake completes.
Invariant 1's negation (an identified-but-not-handshaken connection) is
\emph{reachable} here; Invariant 2 still holds, isolating the Gate 1
failure. The authoritative model is
\texttt{crosshost\_connection\_lifecycle.tex}.}}
\vspace{1em}

% ============================================================
\section{Introduction}
% ============================================================

The same-host Hub-to-Display leg is an \field{AF\_UNIX} socket whose
\texttt{0700} permission \emph{is} the trust model: anything that can open
the path is already a process the same OS user controls.  Cross-host
removes that free lunch --- anyone who can route a packet to the
Display's port can attempt to speak the Hub protocol --- so a
newly-connected peer must clear two gates before its content is honored,
and the accept path must never let a \msg{SceneMessage} slip past either
one.

A connection's life runs through four observable conditions:

\begin{description}
  \item[accepted.] The TCP connection exists; the socket sits in a
    pending-handshake set, deliberately \emph{not} in the reader set, so
    \texttt{poll\_clients} never reads a frame from it
    (\texttt{CrossHostListener.accept\_pending}).
  \item[TLS verified (Gate 1).] \texttt{do\_handshake()} has returned
    successfully --- mutual certificate verification complete.  Only now
    is the fd promoted into the ordinary reader set, where it may exchange
    \msg{ConnectMessage} like any other connection
    (\texttt{CrossHostListener.pump\_ready}).
  \item[identified (Gate 2).] A \msg{ConnectMessage} arrived and its
    self-reported \field{hub\_id} hostname matched the peer certificate's
    SAN (\texttt{CrossHostVerification}); \texttt{handle\_connect}
    recorded the fd as \texttt{kind="hub"}.  A mismatch instead closes the
    connection, never a degraded-trust fallback.
  \item[receiving content.] Only an identified \texttt{kind="hub"} fd has
    its \msg{SceneMessage} processed (\texttt{reject\_scene\_unless\_hub}).
\end{description}

The hazard the design closes is an accept path that hands a socket to the
ordinary reader set --- making its frames deliverable --- before its TLS
handshake, or before its \msg{ConnectMessage}, has completed, letting a
\msg{SceneMessage} land against an unverified or not-yet-identified
connection.  The handshake itself is driven non-blocking, one step per
frame under a bounded deadline, precisely so a stalled peer cannot freeze
the render thread; but the correctness property this model concerns is
the \emph{ordering}, not the timing: content is honored only after both
gates.

Three properties must hold across every interleaving:

\begin{description}
  \item[I1 --- no identify before the handshake (Gate 1).] No connection
    is identified until its TLS handshake has completed.  Because content
    requires identification (I2), this is what keeps a \msg{SceneMessage}
    from an un-handshaken peer out.
  \item[I2 --- no content before identification (Gate 2).] No
    \msg{SceneMessage} is processed against a connection that has not
    been identified (SAN-cross-checked \texttt{kind="hub"}).
  \item[I3 --- deadlock freedom.] Some operation is always enabled.
\end{description}

Together I1 and I2 are exactly Invariant 1: content only after \emph{both}
transport verification and \field{hub\_id} cross-check.

% ============================================================
\section{Basic Types}
% ============================================================

\begin{zed}
  [CONN]
\end{zed}

\texttt{CONN} is the set of cross-host connections (file descriptors) the
Display can hold at once; two suffice to exhibit one connection racing
through the gates while another is at a different stage, and three give
ProB room to interleave a stalled, a verified, and an identified
connection.  The carrier is bounded by ProB's \texttt{DEFAULT\_SETSIZE}.

% ============================================================
\section{State}
% ============================================================

The state is flat: one schema of five subsets of \texttt{CONN}, each
recording a condition a connection has reached.  A connection may hold
several at once (an identified one is also handshaken and readable); the
subset relations that say so are I1 and I2, and are deliberately
\emph{not} written into the schema.

\begin{schema}{Conn}
  known : \power CONN \\
  readable : \power CONN \\
  handshook : \power CONN \\
  identified : \power CONN \\
  served : \power CONN
\where
  readable \subseteq known \\
  handshook \subseteq known \\
  identified \subseteq known \\
  served \subseteq known
\end{schema}

\texttt{known} is every connection currently open (accepted, not yet
dropped).  \texttt{readable} is the ordinary reader set ---
\texttt{SocketListener}'s client dict, the membership \texttt{poll\_clients}
tests before it ever delivers a frame to a handler; a connection outside
it can have no message of any kind processed.  \texttt{handshook} is
Gate 1 complete.  \texttt{identified} is Gate 2 complete.  \texttt{served}
records every connection against which at least one \msg{SceneMessage} has
been processed --- the observable this model exists to constrain.  The
four containments are structural hygiene: dropping a connection removes it
from every set, so nothing outlives \texttt{known}.  The \emph{ordering}
containments --- \texttt{identified \subseteq handshook} (I1) and
\texttt{served \subseteq identified} (I2) --- are the safety properties,
checked as reachability of their negations (Section~\ref{sec:verify}), not
enforced as guards here.

% ============================================================
\section{Initialization}
% ============================================================

\begin{schema}{Init}
  Conn'
\where
  known' = \emptyset \\
  readable' = \emptyset \\
  handshook' = \emptyset \\
  identified' = \emptyset \\
  served' = \emptyset
\end{schema}

A single initial state: no connection exists --- the cold start before any
peer has connected to the cross-host listener.

% ============================================================
\section{Operations}
% ============================================================

\subsection{Accept --- into the pending set, not the reader set}

A peer completes the TCP handshake.  \texttt{c?} becomes known but is
\emph{not} made readable: it sits in the pending-handshake set, kept
separate from the reader set exactly so no frame can be read from it until
its TLS handshake is verified complete.  This is the ordering the whole
design turns on --- \texttt{accept\_pending} wraps the raw socket
non-blocking with \field{do\_handshake\_on\_connect=False} and holds it
apart from \texttt{\_clients}/\texttt{\_readers}.

\begin{schema}{Accept}
  \Delta Conn \\
  c? : CONN
\where
  c? \notin known \\
  known' = known \cup \{ c? \} \\
  readable' = readable \cup \{ c? \} \\
  handshook' = handshook \\
  identified' = identified \\
  served' = served
\end{schema}

\subsection{HandshakeComplete --- Gate 1, promote to the reader set}

\texttt{do\_handshake()} returns successfully for a pending connection.
Only now is \texttt{c?} both recorded as handshaken and promoted into the
reader set --- \texttt{pump\_ready} returning the socket and the caller
moving it out of pending into \texttt{\_clients}.  The two happen together:
there is no state in which a connection is readable but not handshaken.

\begin{schema}{HandshakeComplete}
  \Delta Conn \\
  c? : CONN
\where
  c? \in known \\
  c? \notin handshook \\
  handshook' = handshook \cup \{ c? \} \\
  readable' = readable \cup \{ c? \} \\
  known' = known \\
  identified' = identified \\
  served' = served
\end{schema}

\subsection{Identify --- Gate 2, SAN match}

A \msg{ConnectMessage} arrives on a readable connection and its
\field{hub\_id} hostname matches the peer certificate's SAN.
\texttt{handle\_connect} records the fd as identified.  The guard is
\texttt{c? \in readable}: the message could only be delivered because the
fd is in the reader set.  It does \emph{not} re-check \texttt{handshook}
--- Gate 1's ordering is what guarantees a readable fd is already
handshaken, and removing that ordering (Section~\ref{sec:fidelity},
variant A) is exactly how a not-yet-verified connection reaches identify.

\begin{schema}{Identify}
  \Delta Conn \\
  c? : CONN
\where
  c? \in readable \\
  c? \notin identified \\
  identified' = identified \cup \{ c? \} \\
  known' = known \\
  readable' = readable \\
  handshook' = handshook \\
  served' = served
\end{schema}

\subsection{SanReject --- Gate 2, SAN mismatch, fail-closed}

The other outcome of Gate 2: the \msg{ConnectMessage}'s \field{hub\_id}
does not match the certificate SAN, so \texttt{handle\_connect} closes the
connection (\texttt{remove\_client}) rather than identifying it --- never a
degraded-trust fallback (Invariant 3).  \texttt{c?} leaves every set.

\begin{schema}{SanReject}
  \Delta Conn \\
  c? : CONN
\where
  c? \in readable \\
  c? \notin identified \\
  known' = known \setminus \{ c? \} \\
  readable' = readable \setminus \{ c? \} \\
  handshook' = handshook \setminus \{ c? \} \\
  identified' = identified \\
  served' = served \setminus \{ c? \}
\end{schema}

\subsection{ProcessScene --- content, only past both gates}

A \msg{SceneMessage} is processed.  The guard is the conjunction of both
gates: \texttt{c? \in readable} (the frame was delivered) \emph{and}
\texttt{c? \in identified} (Gate 2 passed) --- the two facts
\texttt{reject\_scene\_unless\_hub} checks before a scene is installed.
This guard \emph{is} I2; weakening it to \texttt{readable} alone
(Section~\ref{sec:fidelity}, variant B) is the W1 gap where an
unidentified fd's scene is installed.

\begin{schema}{ProcessScene}
  \Delta Conn \\
  c? : CONN
\where
  c? \in readable \\
  c? \in identified \\
  served' = served \cup \{ c? \} \\
  known' = known \\
  readable' = readable \\
  handshook' = handshook \\
  identified' = identified
\end{schema}

\subsection{Drop --- an ordinary disconnect}

The connection's socket closes on its own.  \texttt{c?} leaves every set
--- the existing per-fd reaping, applied identically to a cross-host fd.

\begin{schema}{Drop}
  \Delta Conn \\
  c? : CONN
\where
  c? \in known \\
  known' = known \setminus \{ c? \} \\
  readable' = readable \setminus \{ c? \} \\
  handshook' = handshook \setminus \{ c? \} \\
  identified' = identified \setminus \{ c? \} \\
  served' = served \setminus \{ c? \}
\end{schema}

% ============================================================
\section{Invariants}
\label{sec:inv}
% ============================================================

\paragraph{Invariant 1 --- no identify before the handshake (Gate 1).}
\[
  identified \subseteq handshook.
\]
\texttt{Identify} requires \texttt{c? \in readable}, and \texttt{readable}
is granted only by \texttt{HandshakeComplete}, which grants
\texttt{handshook} in the same step --- so \texttt{readable \subseteq
handshook} holds throughout, and hence so does \texttt{identified
\subseteq handshook}.  \texttt{Accept} adds to \texttt{known} only, never
to \texttt{readable}; that is the ordering.  The check
(Section~\ref{sec:verify}) asks whether the negation --- an identified
connection that is not handshaken --- is reachable.

\paragraph{Invariant 2 --- no content before identification (Gate 2).}
\[
  served \subseteq identified.
\]
\texttt{ProcessScene} requires \texttt{c? \in identified} before recording
it in \texttt{served}, and no other operation grows \texttt{served}.  With
I1, this gives \texttt{served \subseteq identified \subseteq handshook}:
content only past both gates --- Invariant 1 of the design in full.

\paragraph{Invariant 3 --- deadlock freedom.}
\texttt{Accept} is enabled for any \texttt{c? \notin known} whenever one
exists; in the sole state where it has no fresh argument
(\texttt{known = CONN}), \texttt{Drop} is enabled.  So some operation is
always enabled.

% ============================================================
\section{Verification: the goals that must flip}
\label{sec:verify}
% ============================================================

Every command below names \emph{this} buggy file, not the correct spec ---
a reader who runs them against this variant sees the defect, which is the
whole point of a fidelity control.  The correct spec
(\texttt{crosshost\_connection\_lifecycle.tex}) returns the opposite
verdict; its Fidelity section narrates the trace.

Type-checking is a \texttt{fuzz} pass (this variant is well-typed --- the
defect is semantic, not syntactic):
\begin{verbatim}
  fuzz docs/crosshost_connection_lifecycle_no_handshake_gate_buggy.tex
\end{verbatim}

\paragraph{I1's negation must be \emph{found} (Gate 1 bypassed).}
\texttt{Accept} here makes a socket readable before its handshake
(\texttt{readable' = readable \cup \{c?\}}), so \texttt{Identify} fires on
a not-yet-handshaken connection.  Trace: \texttt{Accept}(c);
\texttt{Identify}(c) --- \texttt{c} identified while
\texttt{c \notin handshook}.  An unauthenticated peer honored as a Hub,
the mTLS handshake and SAN check bypassed.
\begin{verbatim}
  # must be FOUND against this variant (NOT found against the correct spec)
  probcli docs/crosshost_connection_lifecycle_no_handshake_gate_buggy.tex \
    -model_check -nodead \
    -goal "card({c | c : CONN & c : identified & not(c : handshook)}) > 0" \
    -p DEFAULT_SETSIZE 3
\end{verbatim}

\paragraph{I2 must stay \emph{not found} (this variant isolates Gate 1).}
\texttt{ProcessScene}'s content guard is untouched here, so no
served-but-not-identified state is reachable --- this variant proves
nothing about Gate 2 and must not be mistaken for a Gate 2
counterexample.
\begin{verbatim}
  # must remain NOT found -- ProcessScene still requires identified here
  probcli docs/crosshost_connection_lifecycle_no_handshake_gate_buggy.tex \
    -model_check -nodead \
    -goal "card({c | c : CONN & c : served & not(c : identified)}) > 0" \
    -p DEFAULT_SETSIZE 3
\end{verbatim}

The single change from the correct spec is \texttt{Accept} granting
\texttt{readable} before the handshake; the Gate 2 content guard is intact,
so I2 holds and only I1 flips.  Deadlock-freedom is unaffected and not the
finding here; the correct spec model-checks it.

\end{document}
